\documentclass{article}
\usepackage{graphicx} % Required for inserting images
\usepackage{fancyhdr}
\usepackage{amsmath}
\usepackage{float}


\begin{document}
\pagestyle{fancy}
\fancyhead[RO]{Mehmet Çağatay Kurmaç}
\fancyhead[LO]{Dynamic VB Example}


\section*{1.}

\textbf{E - Step}:\\
Similar to the previous problem, our variables $x$, $w$ and $v$ are gaussian random variables. Which makes our system linear and Gaussian. In this problem, we have access to entire
measurement dataset, thus, we will perform RTS smoothing to find the posterior distribution of the latent variables $x_i$. It is basically a reverse Kalman filter, which starts
from the last time step and goes back to the first time step. The resulting distribution will be the joint distribution of the latent variables, given measurements and the
parameter $\lambda$. So we will have obtained the posterior distribution of the latent variables $p(x|y,\lambda)$. Thus, E-step will be completed this way.

\textbf{M - Step}:\\

We will express the joint distribution of the latent variables and measurements, given the parameter $\lambda$ as follows:
\[
    \begin{aligned}
        p(y_{1:N},x_{1:N}|\lambda) &= p(y_{N}|x_{1:N},y_{1:N-1},\lambda)*p(x_{1:N}|y_{1:N-1},\lambda)*p(y_{1:N-1}|\lambda)\\
                                   &= p(y_{N}|x_{N},\lambda)*p(x_{N}|x_{N-1},\lambda)*p(x_{1:N-1}|y_{1:N-1},\lambda)*p(y_{1:N-1}|\lambda)\\
    \end{aligned}
\]

To obtain the expression in the second line, we used the Markov property of the system. If $X_{k-1}$ is known, then $X_k$ is independent of all previous states and measurements. 
Also, measurements are independent from each other and only depend on the current state. If we continue to expand the joint distribution, we can write the following:
\[
    \begin{aligned}
        p(y_{1:N},x_{1:N}|\lambda) &= p(y_{N}|x_{N},\lambda)*p(x_{N}|x_{N-1},\lambda)*p(x_{N-1}|x_{N-2},\lambda) \cdots p(y_{1}|x_{1},\lambda)*p(x_{1}|\lambda)\\
    \end{aligned}
\]

Then we can obtain the closed form of the joint distribution as the following: 

\[
    p(y_{1:N},x_{1:N}|\lambda) = \prod_{i=1}^{N} p(y_i|x_i,\lambda)*\prod_{i=2}^{N} p(x_i|x_{i-1},\lambda)*p(x_1|\lambda)
\]

For the M-step, we will maximize the expected log-likelihood of the joint distribution with respect to $\lambda$. The expected log-likelihood is given as follows:
\[
    \begin{aligned}
        E_{p(x|y,\lambda)}\left[ ln(p(y_{1:N},x_{1:N}|\lambda)) \right] &= E_{p(x|y,\lambda)}\left[ ln(\prod_{i=1}^{N} p(y_i|x_i,\lambda)*\prod_{i=2}^{N} p(x_i|x_{i-1},\lambda)*p(x_1|\lambda)) \right] \\
        &= E_{p(x|y,\lambda)}\left[ \sum_{i=1}^{N} ln(p(y_i|x_i,\lambda)) + \sum_{i=2}^{N} ln(p(x_i|x_{i-1},\lambda)) + ln(p(x_1|\lambda)) \right] \\
    \end{aligned}
\]

the second sum and the third term are independent of $\lambda$, thus, we can remove them from the expectation. We can simplify our expression as follows:
\[
    \begin{aligned}
            E_{p(x|y,\lambda)}\left[ ln(p(y_{1:N},x_{1:N}|\lambda)) \right] &= E_{p(x|y,\lambda)}\left[ \sum_{i=1}^{N} ln(p(y_i|x_i,\lambda)) \right] + C \\
                                                                            &= \sum_{i=1}^{N} \left[ \frac{1}{2}ln(\lambda ) - \frac{1}{2} ln(2\pi) - \frac{\lambda}{2}E_{p(x|y,\lambda)}\left[(y_i - x_i)^2\right] \right]
    \end{aligned}
\]

The indices of the last term is calculated in E-step with the RTS smoother. Then we can take the derivative of the expected log-likelihood with respect to $\lambda$ and set
it to zero to find the optimal $\lambda$ as follows:
\[
    \begin{aligned}
        \frac{d}{d\lambda} E_{p(x|y,\lambda)}\left[ ln(p(y_{1:N},x_{1:N}|\lambda)) \right] &= \frac{N}{2\lambda} - \frac{1}{2}\sum_{i=1}^{N}E_{p(x|y,\lambda)}\left[(y_i - x_i)^2\right] = 0 \\
        \Rightarrow \lambda &= \frac{N}{\sum_{i=1}^{N}E_{p(x|y,\lambda)}\left[(y_i - x_i)^2\right]}
    \end{aligned}
\]
\newpage
\section*{2.}
We did the formulation in the previous problem. Now, we can apply the framework as given before.
Let's apply this to our problem. Distribution of $\lambda$ is given as $\lambda \sim Gam(\lambda|a,b)$ again. The factorization of the latent variable and parameters is 
given as $p(x_{1:N},\lambda|y_{1:N}) = q_x(x_{1:N})q_\lambda(\lambda)$. As we have explained in the EM solution, $q(x_{1:N})$ is a joint distribution of the latent variables
, given the measurements. Let's find the optimal distribution. 
\[
    \begin{aligned}
        ln(q_x^*(x_{1:N}))  &= E_{\lambda}\left[ \ln p(y,x,\lambda) \right] \\
                            &= E_{\lambda}\left[ \ln p(x_{1:N},y_{1:N}|\lambda) + \ln p(\lambda) \right] \\
                            &= E_{\lambda}\left[ \sum_{i=1}^{N} \ln p(y_i|x_i,\lambda) + \sum_{i=2}^{N} \ln p(x_i|x_{i-1}) + p(x_{1}) + \ln p(\lambda) \right] \\
                            &= E_{\lambda}\left[ \sum_{i=1}^{N} \ln p(y_i|x_i,\lambda) \right] + \sum_{i=2}^{N} \ln p(x_i|x_{i-1}) + p(x_{1}) + C\\
                            &= E_{\lambda}\left[ \sum_{i=1}^{N} \frac{1}{2}ln(\lambda) - \frac{1}{2}ln(2\pi) - \frac{\lambda}{2}(y_i - x_i)^2 \right] + \sum_{i=2}^{N} \ln p(x_i|x_{i-1}) + p(x_{1}) + C\\
                            &= \sum_{i=1}^{N} \frac{1}{2}E_{\lambda}\left[ ln(\lambda) \right] - \frac{1}{2}ln(2\pi) - \frac{1}{2}E_{\lambda}\left[ \lambda \right](y_i - x_i)^2 + \sum_{i=2}^{N} \ln p(x_i|x_{i-1}) + p(x_{1}) + C\\
                            &= \sum_{i=1}^{N} - \frac{1}{2}\frac{a}{b}(y_i - x_i)^2 + \sum_{i=2}^{N} \ln p(x_i|x_{i-1}) + p(x_{1}) + C\\
                            &= - \frac{1}{2}\frac{a}{b}\sum_{i=1}^{N}(y_i - x_i)^2 + \sum_{i=2}^{N} - \frac{1}{2}\ln(2\pi*\sigma_w^2) - \frac{1}{2*\sigma_w^2}(x_i - x_{i-1})^2 + p(x_{1}) + C\\
                            &= - \frac{1}{2}\frac{a}{b}\sum_{i=1}^{N}(y_i - x_i)^2 + \sum_{i=2}^{N} - \frac{1}{2*\sigma_w^2}(x_i - x_{i-1})^2 - \frac{1}{2*\Sigma_{1}^2}(x_1 - \bar{x_{1}})^2 + C\\
    \end{aligned}
\]

The last line is the log of product of Gaussian distributions, which are familiar.
\[
    \ln(q_x^*(x_{1:N})) = \sum_{i=1}^{N} \ln\mathcal{N}(y_i ; x_i, \frac{b}{a}) + \sum_{i=2}^{N} \ln\mathcal{N}(x_i ; x_{i-1}, \sigma_w^2) + \ln\mathcal{N}(x_1 ; \bar{x_1}, \Sigma_1^2) + C
\]

We know this from our first factorization of $p(x_{1:N},y_{1:N}|\lambda)$. This is the joint distribution of latent variables, given the measurements. To obtain the posterior,
namely, optimal distribution of $q_x^*(x_{1:N})$, we can use the RTS smoother. Only difference is the variance of the measurement noise, which is $\frac{b}{a}$ instead of
$\frac{1}{\lambda}$. (Basically the same thing but we will optimize with respect to $a$ and $b$ in the iterations.)

Let's also find the optimal distribution of $\lambda$:

\[
    \begin{aligned}
        ln(q_\lambda^*(\lambda)) &= E_{x}\left[ \ln p(y,x,\lambda) \right] \\
                        &= E_{x}\left[ \ln p(x_{1:N},y_{1:N}|\lambda) + \ln p(\lambda) \right] \\
                        &= E_{x}\left[ \sum_{i=1}^{N} \ln p(y_i|x_i,\lambda) + \sum_{i=2}^{N} \ln p(x_i|x_{i-1}) + p(x_{1}) + \ln p(\lambda) \right] \\
                        &= E_{x}\left[ \sum_{i=1}^{N} \ln p(y_i|x_i,\lambda) \right] + ln(p(\lambda)) + C \\
                        &= E_{x}\left[ \sum_{i=1}^{N} \ln \mathcal{N}(y_i;x_i,\frac{1}{\lambda}) \right] + ln(Gam(\lambda|a,b)) + C\\
                        &= E_{x}\left[ \sum_{i=1}^{N} \frac{1}{2}ln(\lambda) - \frac{1}{2}ln(2\pi) - \frac{\lambda}{2}(y_i - x_i)^2 \right] + (a-1)ln(\lambda) - b\lambda - a ln(b) - ln(\Gamma(a)) + C\\
                        &= \left( (a + \frac{N}{2}) - 1 \right)ln(\lambda) - \left( b + \frac{1}{2} E_{x}\left[ \sum_{i=1}^{N}(y_i - x_i)^2\right] \right)\lambda - a ln(b) - ln(\Gamma(a)) + C
    \end{aligned}
\]

The resulting expression is similar to the natural logarithm of a Gamma distribution, which is $(a-1)ln(\lambda) - b\lambda - a ln(b) - ln(\Gamma(a))$. By comparing the coefficients
, we can find the shape and scale parameters of optimal distribution of $\lambda$ as follows:
\[
    \begin{aligned}
        a^* &= a + \frac{N}{2} \\
        b^* &= b + \frac{1}{2} E_{x}\left[ \sum_{i=1}^{N}(y_i - x_i)^2\right] \\
    \end{aligned}
\]  

Then the estimated $\lambda$ is basically the mean of the optimal distribution of $\lambda$, which is $\frac{a^*}{b^*}$.

\section*{3.}

As same as the previous problem, They are equivalent when the posterior distribution of the latent variables is equal to the optimal distribution of the latent variables.

\[ 
    q(x) = p(x|y,\lambda);
\]

Another observation is that the exact match on the prior for VB matters less when sample count N is large. Because the prior will be dominated by the likelihood, and the 
convergence of the VB algorithm will be faster. Thus, the discrepancy between the EM and VB algorithms will be at the beginning, and smaller. 

For the previous single-point observation case, we can reduce the weight of the prior by setting the shape and scale parameters of the Gamma distribution to be very small. If
the ratio of a and b is initiated as the same as the same as the $\lambda$ in the EM algorithm, the results are exactly the same.

\end{document}